\documentclass[10pt,a4paper]{amsart}
\usepackage{amsmath,amssymb,amsthm,graphicx}
% Retained minimal preamble for the Stacks Project chapter more-algebra.tex (More on Algebra).
% The chapter file pinned-original-source.tex carries only \input{preamble} on its first line, so no
% preamble text lives inside the chapter itself.  The Stacks Project uses one single shared preamble for
% every chapter; the snapshot stages it identically as sources/stacks-chow/preamble.tex and
% sources/stacks-etale-cohomology/preamble.tex.  The block below is copied mechanically, line for line,
% from sources/stacks-chow/preamble.tex lines 14-18 (the Xy-pic setup) and 154-217 (the native theorem
% declarations together with the mathematical macros).  Nothing was added, removed or redefined; the
% wrapper class and the amsmath/amssymb/amsthm/graphicx base packages are supplied by the assembler.
% Chapter: more-algebra.tex -- More on Algebra.
% For commutative diagrams we use Xy-pic
\usepackage[all]{xy}

% We use 2cell for 2-commutative diagrams.
\xyoption{2cell}

% Theorem environments.
%
\theoremstyle{plain}
\newtheorem{theorem}[subsection]{Theorem}
\newtheorem{proposition}[subsection]{Proposition}
\newtheorem{lemma}[subsection]{Lemma}

\theoremstyle{definition}
\newtheorem{definition}[subsection]{Definition}
\newtheorem{example}[subsection]{Example}
\newtheorem{exercise}[subsection]{Exercise}
\newtheorem{situation}[subsection]{Situation}

\theoremstyle{remark}
\newtheorem{remark}[subsection]{Remark}
\newtheorem{remarks}[subsection]{Remarks}

\numberwithin{equation}{subsection}

% Macros
%
\def\lim{\mathop{\mathrm{lim}}\nolimits}
\def\colim{\mathop{\mathrm{colim}}\nolimits}
\def\Spec{\mathop{\mathrm{Spec}}}
\def\Hom{\mathop{\mathrm{Hom}}\nolimits}
\def\Ext{\mathop{\mathrm{Ext}}\nolimits}
\def\SheafHom{\mathop{\mathcal{H}\!\mathit{om}}\nolimits}
\def\SheafExt{\mathop{\mathcal{E}\!\mathit{xt}}\nolimits}
\def\Sch{\mathit{Sch}}
\def\Mor{\mathop{\mathrm{Mor}}\nolimits}
\def\Ob{\mathop{\mathrm{Ob}}\nolimits}
\def\Sh{\mathop{\mathit{Sh}}\nolimits}
\def\NL{\mathop{N\!L}\nolimits}
\def\CH{\mathop{\mathrm{CH}}\nolimits}
\def\proetale{{pro\text{-}\acute{e}tale}}
\def\etale{{\acute{e}tale}}
\def\QCoh{\mathit{QCoh}}
\def\Ker{\mathop{\mathrm{Ker}}}
\def\Im{\mathop{\mathrm{Im}}}
\def\Coker{\mathop{\mathrm{Coker}}}
\def\Coim{\mathop{\mathrm{Coim}}}

% Boxtimes
%
\DeclareMathSymbol{\boxtimes}{\mathbin}{AMSa}{"02}

%
% Macros for moduli stacks/spaces
%
\def\QCohstack{\mathcal{QC}\!\mathit{oh}}
\def\Cohstack{\mathcal{C}\!\mathit{oh}}
\def\Spacesstack{\mathcal{S}\!\mathit{paces}}
\def\Quotfunctor{\mathrm{Quot}}
\def\Hilbfunctor{\mathrm{Hilb}}
\def\Curvesstack{\mathcal{C}\!\mathit{urves}}
\def\Polarizedstack{\mathcal{P}\!\mathit{olarized}}
\def\Complexesstack{\mathcal{C}\!\mathit{omplexes}}
% \Pic is the operator that assigns to X its picard group, usage \Pic(X)
% \Picardstack_{X/B} denotes the Picard stack of X over B
% \Picardfunctor_{X/B} denotes the Picard functor of X over B
\def\Pic{\mathop{\mathrm{Pic}}\nolimits}
\def\Picardstack{\mathcal{P}\!\mathit{ic}}
\def\Picardfunctor{\mathrm{Pic}}
\def\Deformationcategory{\mathcal{D}\!\mathit{ef}}

% --- appended 2026-10-12 (Chief): the project own definitions of the reference/slogan/history
% display environments, copied mechanically from the shared Stacks preamble lines 8-12
% (\usepackage{verbatim} plus the three \newenvironment lines, which the chapter preamble keeps
% verbatim from the project preamble). Nothing added, removed or redefined: these environments
% carry no mathematical content and the project itself typesets them as comments. Without these
% three lines a counted statement whose own text uses them fails with Environment slogan undefined.
\usepackage{verbatim}
\newenvironment{reference}{\comment}{\endcomment}
\newenvironment{slogan}{\comment}{\endcomment}
\newenvironment{history}{\comment}{\endcomment}
\usepackage{cleveref}
\crefname{originaldefinition}{original definition}{original definitions}
\Crefname{originaldefinition}{Original definition}{Original definitions}
\makeatletter
\newlabel{definition-projective-dimension}{{0A5N}{source}{}{sourcepointer-1}{}}
\newlabel{definition-projective-dimension@cref}{{[originaldefinition][0][]0A5N}{[1][0][]source}{}{}{}}
\makeatother
\crefname{originallemma}{original lemma}{original lemmas}
\Crefname{originallemma}{Original lemma}{Original lemmas}
\makeatletter
\newlabel{lemma-perfect}{{0658}{source}{}{sourcepointer-2}{}}
\newlabel{lemma-perfect@cref}{{[originallemma][0][]0658}{[1][0][]source}{}{}{}}
\makeatother
\crefname{originallemma}{original lemma}{original lemmas}
\Crefname{originallemma}{Original lemma}{Original lemmas}
\makeatletter
\newlabel{derived-lemma-morphisms-from-projective-complex}{{064B}{source}{}{sourcepointer-3}{}}
\newlabel{derived-lemma-morphisms-from-projective-complex@cref}{{[originallemma][0][]064B}{[1][0][]source}{}{}{}}
\makeatother
\crefname{originallemma}{original lemma}{original lemmas}
\Crefname{originallemma}{Original lemma}{Original lemmas}
\makeatletter
\newlabel{derived-lemma-split}{{05QT}{source}{}{sourcepointer-4}{}}
\newlabel{derived-lemma-split@cref}{{[originallemma][0][]05QT}{[1][0][]source}{}{}{}}
\makeatother
\title{Vanishing of Hom into a complex supported below the projective amplitude, and the resulting unique splitting}
\author{The Stacks Project Authors}
\date{stacks-project:more-algebra.tex:local-snapshot-20261012; GNU FDL 1.2 or later}
\begin{document}
\maketitle
\setcounter{subsection}{1}
\noindent\textit{Editorial note.} All delivered mathematical text is the author's own, copied from the pinned Stacks Project chapter file more-algebra.tex; no statement or proof step was rewritten, reordered or removed. The counted statement is the labeled lemma environment at source lines 19703-19718 (canonical tag 0BCG) and the counted proof is the proof environment at source lines 19720-19741, which immediately follows it and is its own proof.\par\medskip
\noindent\textit{Editorial note.} The four references are two to results of this chapter and two to results of another chapter; the kinds recorded in the stubs are the type words appearing in the labels themselves, as the lane's convention requires, not the names of the environments in which the labels live.\par\medskip
\noindent\textit{Editorial note.} The preamble span runs to line 90, the current end of the retained preamble, which includes the appended environment definitions; the counted text does not use them, and the span is taken whole so that the file version is unambiguous.\par\medskip
\noindent\textit{Editorial note.} The source URL is the source-level Stacks Project URL, as required when more than one card is produced from one source; the per-claim tag is carried by claim\_key and by this note. The counted statement and proof cite no bibliography entry, so this card needs no bibliography.\par\medskip
\noindent\textit{Allowed standard prerequisites (editorial).}\begin{itemize}
\item An object of projective amplitude in an interval is represented by a complex of projective modules vanishing outside that interval, and a perfect object of tor amplitude there is represented by finite projective modules. invoked at proof lines 19720-19727. Provides the representative of the first object in the vanishing computation, in both the general and the perfect case.
\item An object whose cohomology vanishes above a bound is quasi-isomorphic to its truncation below that bound, hence representable by a complex vanishing there. used at proof lines 19727-19731. Puts the second object in the form in which no nonzero map of complexes into it can exist.
\item Maps in the derived category from a bounded complex of projectives are maps of complexes, and a distinguished triangle whose first Hom group vanishes splits. used at proof lines 19731-19739. Gives assertion (1), then the splitting of the triangle in assertion (2) and the uniqueness in assertion (3).
\end{itemize}\medskip
\ifdefined\hypertarget\hypertarget{sourcepointer-1}{}\fi
\noindent\textit{Source reference.} Original definition 0A5N: sources/stacks-more-algebra/pinned-original-source.tex, source lines 17394--17394. The reference is cited at proof lines 19721-19722 to justify representing the object by a complex of projective modules vanishing outside the interval; the rest of the proof is inside the counted span.\par\medskip
\ifdefined\hypertarget\hypertarget{sourcepointer-2}{}\fi
\noindent\textit{Source reference.} Original lemma 0658: sources/stacks-more-algebra/pinned-original-source.tex, source lines 18889--18889. The reference handles the alternative hypothesis that the object is perfect of tor amplitude in the interval, cited at proof lines 19724-19727; it is a pointer and the computation that follows is inside the counted span.\par\medskip
\ifdefined\hypertarget\hypertarget{sourcepointer-3}{}\fi
\noindent\textit{Source reference.} Original lemma 064B: sources/stacks-chow/tags.txt, source lines 7860--7860. The reference identifies the Hom group of the derived category with the group of maps of complexes from a bounded complex of projectives, used at proof lines 19732-19734 to conclude the vanishing; the disjointness of supports that makes it apply is inside the counted span.\par\medskip
\ifdefined\hypertarget\hypertarget{sourcepointer-4}{}\fi
\noindent\textit{Source reference.} Original lemma 05QT: sources/stacks-chow/tags.txt, source lines 7387--7387. The reference is the criterion saying that a distinguished triangle splits when the relevant Hom group vanishes, applied at proof lines 19736-19738 to get assertion (2); the vanishing it consumes is established inside the counted span.\par\medskip
\begin{lemma}
\label{lemma-splitting-unique}
Let $R$ be a ring. Let $K$ and $L$ be objects of $D(R)$.
Assume $L$ has projective-amplitude in $[a, b]$, for example if $L$
is perfect of tor-amplitude in $[a, b]$.
\begin{enumerate}
\item If $H^i(K) = 0$ for $i \geq a$, then
$\Hom_{D(R)}(L, K) = 0$.
\item If $H^i(K) = 0$ for $i \geq a + 1$, then given any distinguished
triangle $K \to M \to L \to K[1]$
there is an isomorphism $M \cong K \oplus L$
in $D(R)$ compatible with the maps in the distinguished triangle.
\item If $H^i(K) = 0$ for $i \geq a$, then the isomorphism
in (2) exists and is unique.
\end{enumerate}
\end{lemma}

\begin{proof}
The assumption that $L$ has projective-amplitude in $[a, b]$ means we
can represent $L$ by a complex $L^\bullet$ of projective $R$-modules
with $L^i = 0$ for $i \not \in [a, b]$, see Definition
\ref{definition-projective-dimension}.
If $L$ is perfect of tor-amplitude in $[a, b]$, then
we can represent $L$ by a complex $L^\bullet$ of finite projective $R$-modules
with $L^i = 0$ for $i \not \in [a, b]$, see Lemma \ref{lemma-perfect}.
If $H^i(K) = 0$ for $i \geq a$, then $K$ is quasi-isomorphic
to $\tau_{\leq a - 1}K$. Hence we can represent $K$ by a complex
$K^\bullet$ of $R$-modules with $K^i = 0$ for $i \geq a$. Then we obtain
$$
\Hom_{D(R)}(L, K) = \Hom_{K(R)}(L^\bullet, K^\bullet) = 0
$$
by Derived Categories, Lemma
\ref{derived-lemma-morphisms-from-projective-complex}.
This proves (1). Under the hypotheses of (2) we see that
$\Hom_{D(R)}(L, K[1]) = 0$ by (1), hence
the distinguished triangle is split by
Derived Categories, Lemma \ref{derived-lemma-split}.
The uniqueness of (3) follows from (1).
\end{proof}

\end{document}
